%%%%%%%%%%%%
%
% $Autor: Wings $
% $Datum: 2019-03-05 08:03:15Z $
% $Pfad: SetUp.tex $
% $Version: 4250 $
% !TeX spellcheck = en_GB/de_DE
% !TeX encoding = utf8
% !TeX root = manual 
% !TeX TXS-program:bibliography = txs:///biber
%
%%%%%%%%%%%%

\chapter{Set Up}

\section{System Setup}

This section describes the step-by-step setup procedure for assembling and preparing the \textbf{License Plate Detection System} for operation.  
It covers the physical arrangement, connections, and software initialization on the NVIDIA Jetson Nano platform.

\vspace{0.4cm}
\noindent
\textbf{Step 1: Hardware Placement and Assembly}
\begin{itemize}
	\item Place the \textbf{Jetson Nano (reComputer J1010)} unit on a flat, well-ventilated surface to ensure proper passive cooling.  
	\item Mount the \textbf{Microsoft LifeCam Show} securely on the \textbf{vertical metallic stand} and position it to face the target area (e.g., parking lane or vehicle entry point).  
	\item Connect the camera to one of the Jetson Nano’s USB 2.0 ports.  
	\item Connect the HDMI display cable from the Nano to the external monitor.  
	\item Attach the power adapter (5V/4A USB-C) but do not switch on the power yet.
\end{itemize}

\vspace{0.3cm}
\noindent
\textbf{Step 2: Storage and Operating System Preparation}
\begin{itemize}
	\item Insert the microSD card preloaded with \textbf{JetPack 5.1.2 (Ubuntu 20.04)} into the Jetson Nano’s microSD slot.  
	\item Ensure that all connections (camera, HDMI, keyboard/mouse, and power) are properly seated.  
	\item Power on the system to boot into the Ubuntu desktop environment.  
	\item Log in to the user account and open the terminal.
\end{itemize}

\vspace{0.3cm}
\noindent
\textbf{Step 3: Software Initialization and Testing}
\begin{itemize}
	\item Open a terminal and navigate to the project directory:  
	\texttt{cd ~/LicensePlateDetection/}
	\item Run the detection script:  
	\texttt{python3 main.py}
	\item The system will start the \textbf{YOLOv8} model, initialize the camera feed, and display detection results on-screen.  
	\item Confirm that the detection output and recognized text are being logged into the \texttt{plates.db} SQLite database.
\end{itemize}

\vspace{0.5cm}
\noindent
\textbf{Step 4: Verification and Adjustment}
\begin{itemize}
	\item Adjust the camera tilt and distance until the license plate region appears clearly within the detection frame.  
	\item Verify focus and illumination to ensure consistent recognition accuracy.  
	\item Check that power and HDMI connections remain stable during prolonged operation.
\end{itemize}

\vspace{0.6cm}
\noindent
\textbf{System Setup Diagram:}

\begin{figure}[h!]
	\centering
	\resizebox{1.2\textwidth}{!}{%
		\begin{tikzpicture}[
			node distance=2.4cm, 
			>=stealth,
			font=\large,
			every node/.style={align=center}
			]
			
			% Define styles
			\tikzstyle{block} = [rectangle, draw=blue!70!black, fill=blue!10, 
			rounded corners, minimum height=1cm, minimum width=3.3cm, text centered, font=\small]
			\tikzstyle{process} = [rectangle, draw=green!60!black, fill=green!10, 
			rounded corners, minimum height=1cm, minimum width=3.6cm, text centered, font=\small]
			\tikzstyle{db} = [cylinder, shape border rotate=90, draw=orange!90!black, fill=orange!10, 
			minimum height=1.3cm, minimum width=1.1cm, aspect=0.4, font=\small]
			\tikzstyle{arrow} = [->, thick]
			
			% --- LAYER 1: Hardware devices ---
			\node[block] (camera) {\faCamera \\ \textbf{Microsoft LifeCam Show}};
			\node[block, right=4.0cm of camera] (nano) {\faMicrochip \\ \textbf{NVIDIA Jetson Nano}\\(reComputer J1010)};
			\node[block, right=4.0cm of nano] (display) {\faDesktop \\ \textbf{HDMI Display}};
			\node[block, below=1.7cm of nano, fill=gray!10, draw=gray!60!black] (power) {\faBolt \\ 5V / 4A USB-C Power Supply};
			
			% --- LAYER 2: Software / Processing pipeline ---
			\node[process, below=2cm of power, xshift=-3.9cm] (os) {\textbf{JetPack 5.1.2 (Ubuntu 20.04)}};
			\node[process, right=4.0cm of os] (yolo) {\textbf{YOLOv8 Inference Engine}};
			\node[process, right=4.0cm of yolo] (ocr) {\textbf{EasyOCR Recognition}};
			\node[db, below=2cm of yolo, yshift=-0.2cm] (database) {\textbf{SQLite}\\plates.db};
			
			% --- ARROWS: Hardware layer ---
			\draw[arrow] (camera) -- node[above, yshift=0.2cm, text width=3.5cm]{\scriptsize Video Stream \\ (USB 2.0)} (nano);
			\draw[arrow] (nano) -- node[above, yshift=0.2cm, text width=3.5cm]{\scriptsize Annotated Output \\ (HDMI)} (display);
			\draw[arrow] (power) -- node[right]{\scriptsize Power Input (USB-C)} (nano);
			
			% --- ARROWS: Software layer ---
			\draw[arrow, dashed, blue!80!black] (os) -- node[above, yshift=0.15cm, text width=3cm, align=center]{\scriptsize Boot + Runtime Environment} (yolo);
			\draw[arrow, dashed, blue!80!black] (yolo) -- node[above, yshift=0.15cm, text width=3cm, align=center]{\scriptsize Detected Plate Region} (ocr);
			\draw[arrow, dashed, blue!80!black] (ocr) -- node[above, yshift=0.15cm, text width=3cm, align=center]{\scriptsize Extracted Text + Confidence} (database);
			
			% --- CONNECTIONS BETWEEN LAYERS ---
			\draw[arrow, thick, gray!70!black, bend right=20] (nano.south) to node[left, xshift=-0.2cm]{\scriptsize Runs OS + Models} (os.north);
			\draw[arrow, thick, orange!80!black, bend left=20] (database.north) to node[right, xshift=0.2cm]{\scriptsize Data Logging + Query} (display.south);
			
			% --- DECORATIVE BACKGROUND BOXES (grouping) ---
			\node[draw=blue!40!black, thick, rounded corners, fit=(camera)(nano)(display)(power), label={[yshift=0.2cm]above:\textbf{Hardware Layer}}] {};
			\node[draw=green!40!black, thick, rounded corners, fit=(os)(yolo)(ocr)(database), label={[yshift=-0.3cm]below:\textbf{Software and Data Flow Layer}}] {};
			
		\end{tikzpicture}%
	}
	\caption{Comprehensive System Setup Flow – Hardware and Software Integration}
	\label{fig:setup_complex_tikz}
\end{figure}

\vspace{0.4cm}
\noindent
The setup diagram (Figure~\ref{fig:setup_complex_tikz}) visually represents the physical connection flow.  
All communication occurs locally between hardware components, ensuring low latency and complete offline operation.  
This modular arrangement allows easy relocation and quick system reassembly at different deployment sites.